\section{高级功能：Extended Thinking 与 MCP}

\subsection{Extended Thinking}

对于复杂问题，可以开启 Extended Thinking 模式让 Claude 在回答前花更多时间推理。

\begin{importantbox}{何时使用 Extended Thinking}
\textbf{适合}：复杂架构决策、棘手的多因素 debug、多步骤重构方案、多方案利弊权衡。\\
\textbf{不需要}：简单代码修改、格式化/重命名、已经很明确的 bug 修复。\\
\textbf{小技巧}：在提示词里写 \texttt{ultrathink} 可以强制触发最高推理深度。
\end{importantbox}

\subsection{MCP Server}

MCP（Model Context Protocol）让 Claude Code 能和外部工具通信——GitHub、数据库、Slack、Notion 等。配置在项目根目录的 \texttt{.mcp.json} 里。

\begin{warningbox}{MCP 的上下文开销}
每个 MCP server 会额外消耗 100--500+ token 的工具定义。用 \texttt{/mcp} 命令查看每个 server 的开销，禁用不用的 server。\textbf{不要贪多}——常用的 2--3 个就够了，装太多反而影响代码分析质量。
\end{warningbox}

\subsection{IDE 集成}

Claude Code 不仅可以在终端使用，还有 VS Code 扩展和 JetBrains 插件：

\begin{itemize}[leftmargin=1.5em]
  \item \textbf{VS Code}：并排 diff 视图、\texttt{Option+K} 快速引用文件、多标签对话、\texttt{Cmd+Shift+Esc} 新建对话
  \item \textbf{JetBrains}：IntelliJ/PyCharm/WebStorm 等均支持，在 Plugins 市场搜索安装
\end{itemize}

\subsection{本章小结}

Extended Thinking 适合复杂推理任务。MCP 扩展了 Claude 的能力边界，但要注意上下文开销。IDE 集成让工作流更顺滑。

%% ====================================================================
